\section{周王室之末：当命令再也压不住声音}

\subsection{厉王：资源、边防与朝廷里的裂纹}

厉王进入晚西周王序时，王室仍能发号施令，也仍须维持边防、赏赐和朝廷中贵族的协调。正因为这些事务都要动用资源，王命的每一次伸展都可能触到既有的利益和名分。现存材料不足以让我们逐项复原王室的收入、支出或具体政策；能够看见的，是后来的危机叙事把财政、用人与朝廷关系的紧张集中到了厉王一朝。

《国语·周语上》说厉王“专利”，《史记》又以相近的线索铺陈其失政。这里的“利”究竟指山林川泽等资源、某种收益权，还是后世作者用以概括王室与贵族、国人冲突的说法，至今不能断定。它不能被径直译作现代意义的“垄断”，更不能据此列出一套确证的税制。不过，边防调发、赏赐维系与朝廷竞争都需要可支配的物资；后出叙事之所以反复把冲突写到资源上，至少保留了一个问题：当王室试图更直接地掌握收益时，谁会失去原有的位置，谁又能迫使它退让？

\begin{debate}[“专利”能说明什么]
《国语》与《史记》保存的是回望西周末年的政治叙事，不是厉王朝的财政账簿。它们可以用来追问晚西周的资源压力、精英关系和王权边界；不能用来断言某一项政策的名称、实施范围和收益数额。把道德化的“暴君”图像还原成制度史，第一步正是承认材料在这里的限度。
\end{debate}

厉王危机也不能只写成朝廷里的一场争吵。中心对资源的要求、边防所需的兵粮、贵族集团可动员的人力，以及王室命令获得回应的程度，原本相互支撑；任何一环变弱，另外几环的成本都会升高。从现存材料看，厉王成为这些压力的集中人物，并不等于压力始于他一人。后来的出奔，才使潜在的裂纹变成了公开的政局中断。

\subsection{共和元年：一条年表从危机里开始}

\eventanchor{WZ-0841-GONGHE}{西周·共和}

中国传统连续纪年从前841年开始。这个起点并不安静：它不是一位新君登基的吉日，而是厉王出奔之后，王室不得不用非常安排维持局面的年份。史书把它称为“共和”。从此以后，年代可以一岁一岁地往下数；在这之前，许多王年的绝对位置仍像隔着雾。

《国语·周语上》写厉王“专利”，又写他试图压住批评。故事里，召公对厉王说：

\begin{textwindow}[声音与河水]
“防民之口，甚于防川。川壅而溃，伤人必多；民亦如之。”

出处为《国语·周语上》。它把召公的进谏置于厉王压制谤议的危机中，是一篇借往事劝戒君主的国别史叙事；这段警语能说明后世如何把冲突失控理解为西周政治危机，不能当作共和前夕的逐字实录。
\end{textwindow}

这段话后来常被读成一句关于言论的格言。放回西周末年，它更像一次政治警告：王室需要山林、川泽和其他资源来维持军政与赏赐；贵族、国人和王室围绕这些资源的关系已经紧张。若把反对者一概当作需要消除的噪音，原本还能谈、还能调停的冲突，就可能汇成无法收拾的力量。

厉王出奔以后，祭祀、军政与王统不能长久无人承接。传世叙事一说周、召二相共同执政，另一条线索则说共伯和一度“干王位”。两说都在解释同一个非常事实：王室没有按通常的父子继承立即恢复一位可无争议行使王权的君主。共和的起止年可以排列，谁实际以何种名义主持其间政务，却不能写成已经解决的问题。

这种非常安排使秩序没有立刻断裂，也显示出王室尚有可调用的精英网络；同时，它把王权受制的情形暴露在外。共和不是西周立刻结束的时刻，却让中心和核心支持者之间被磨薄的信任变得清晰可见。

\begin{debate}[“共和”是谁在执政]
《史记》说“周召共和”；《竹书纪年》则有“共伯和干王位”的表述。近出材料使争论没有完全终结。对这一章而言，关键不在把两说强分高下，而在于：厉王出奔后，王室需要一个不能按常规继承解释的政治安排，才能让秩序不断裂。
\end{debate}

\subsection{宣王复位：恢复命令，并不等于恢复旧日条件}

前828年，宣王即位，共和结束，正式继承重新接上王统。后世常以“中兴”概括这一转折；这个词首先提醒我们，恢复王位和重建朝廷命令是当时必须完成的事，而不是一张可以证明全面复原的成绩单。关于其初年具体措施，可靠材料有限；毛公鼎等晚西周金文以及传世王年材料，只能让研究者继续讨论朝廷命令、贵族关系和晚西周的年代位置。

宣王的即位说明非常行政不能永久替代王位继承，但它没有取消先前已经显露的难题。王室若要维持边防与朝廷，仍须取得兵力、物资和精英合作；诸侯与贵族也会据此衡量王命的代价和回报。因而，“中兴”应被读作一项复位的政治工程，而非倒推幽王覆亡之前一切危机均已消失。宣王之后，幽王继承王位，材料对于其初年尤其稀少；正因如此，不能把最后的崩溃简单归为一朝骤变。

\subsection{幽王：继承的私事成为联盟的断点}

幽王即位时，西周还没有倒下。王室仍有宗周，仍有军政与礼仪的名分，也仍依靠封国和贵族网络。可这些东西已不能像从前那样彼此咬合。边防需要兵，兵需要粮；王室需要诸侯，诸侯也在计算王室还能给什么、还能要求什么。幽王初年的材料不足，不能把这一结构压力逐年摊开；到继承冲突公开时，它已为宫廷决定外溢为封国、外戚和边疆力量的对抗提供了条件。

较早与后出的材料都留下了废申后、废太子宜臼、改立褒姒之子伯服的叙事线索。《国语·郑语》《史记·周本纪》以及《左传》所见材料并不允许复原完整程序或每个人的动机，但继承危机的轮廓可信：嫡庶名分和婚姻联盟被重新排序，申侯与幽王的关系随之破裂；宜臼又在申地或诸侯支持下保有竞争王位的政治位置。原本属于宫廷的废立，在宗法和封国网络中从来不是私事。

\eventanchor{WZ-0771-HAOJING}{西周·镐京陷落}

前771年，幽王死，镐京陷落。后世最爱讲“烽火戏诸侯”：褒姒不笑，幽王便点燃烽火，诸侯白跑一场，及至真正有难便不再来援。这故事太完整了，完整得像专为一个王朝的覆亡准备。它也恰恰因此值得怀疑。

较早的材料更接近一场继承危机和政治联盟的崩裂。申侯与犬戎等力量如何结盟、犬戎究竟是怎样的政治和军事集团、行动路线又为何，均有争议；可以较有把握地说，内部裂解使西部的防御体系失去协调，边疆力量得以进入关中危机。王位继承牵动了外戚、诸侯、边疆军队和关中的防线，军事资源并未仍由镐京单独控制。西周最后一次危机，不是一个人荒唐到失去天下，而是一个本已松动的体系，终于在继承问题上同时断开。

\subsection{镐京之后：王统东移，王畿不再随行}

镐京遭劫之后，宗周的都城体系已难恢复。丰镐晚期层位和传世史书都支持西部中心失去运作条件这一结论，但宫室、仓储和聚落受破坏的范围不能由此推成一幅整齐的废墟图景。关中并非一夜无人，周原及其周边聚落在失去王都环境后如何调整居住与生产，也只能在区域调查和分期材料中看见趋势，不能逐村追踪去向。

平王转向洛邑，是在镐京无法恢复、成周仍有宫室、祭祀和交通节点可供承接的条件下作出的选择。郑氏家族护卫东迁、晋文侯辅佐平王的记忆，分别见于郑、晋相关的后世谱系叙事；它们提示新朝廷需要诸侯和卿族的支持，但人物的先后、封赏和行动规模并不清楚。东迁的可靠性不依赖这些细节：王统在洛邑得以续接，而重建朝会关系本身已显示它不能再完全倚赖旧有的自有军力。

周还在；祭祀、册命和“天子”的称谓也还在。但它所依靠的关中根基、直接军力和资源中心已经失去。成周可以承接名分，却远不能等同于宗周的物质与军事基础。此后的周王室仍能参与界定合法性，却越来越依赖诸侯来维持这种合法性能够发挥的场合。

\begin{causalchain}[从共和到东迁]
\eventref{WZ-0841-GONGHE}{西周·共和}让王室权威的裂缝显形；\eventref{WZ-0771-HAOJING}{西周·镐京陷落}则让继承危机、封国联盟和边防失效在同一时刻汇合。东迁不是一场搬家，而是统治能力的重心转移：周王室保住了名分，却失去了支配名分的物质底座。
\end{causalchain}

\begin{sourcenote}[传说、王年与政治过程]
厉王段主要依靠《国语·周语上》与《史记》，并须与《竹书纪年》、史墙盘所见王序相互辨读；“专利”的政策细节不能确证。共和的前841年至前828年是可靠年代锚点，但周、召二相与共伯和两说不能强作一解。宣王的复位见《史记》《竹书纪年》及晚西周金文研究，所谓“中兴”不宜扩大为全面复原。

幽王段可参看《国语》《左传》《史记》及西周晚期《诗经·小雅》所呈现的危机氛围。烽火戏诸侯的完整情节主要见于《史记》，无同时代证据；它不应替代对继承、军政、边防和联盟的分析。前771年镐京陷落与前770年前后东迁的年代脊柱较稳，而申侯、犬戎、郑、晋在危机中的具体行动仍须区分传世叙事、地方谱系记忆、金文考古和现代重建。史书的文学性在这里尤其强：它能把复杂制度崩解凝成一个人、一位美人、一团烽火；读者应当记住传说的力量，也应回到结构压力与材料限度中去。
\end{sourcenote}
